Persistent component motives can change conversational behavior without improving task accuracy. We characterize a shared-backbone system with generic adversaries, biography-plus-motive components, and motive-only components. Across 150 paired synthetic planning families, three generator seeds, and 10,800 intervention opportunities per condition, biography-plus-motive components make 477 direct edits versus 5,408 for generic adversaries, while using 4,618 recollection actions versus 828. Typed preference repetition is high in both character and motive-only conditions, but API labels do not establish the meaning of a claim. We evaluate two frozen model-family assessors on 80 held-out synthetic controls each, 612 archive-sample records, 32 fresh controlled events, and stability probes. Phi passes 0 of seven declared screening dimensions and Qwen passes 3; semantic outputs remain separately reported, with calibration and natural-dialogue validity distinguished. A complete 64-conversation experiment changes compatible to incompatible same-field preferences. Typed matching changes sharply, but most matching follows an existing position; paired task success is 5/16 versus 7/16 families, with wide uncertainty. Antagonist-free C5-minus-generic accuracy differences are -3.89 and -4.17 percentage points for Qwen and Phi, respectively. Traceable dialogues support a bounded account of advocacy, evidence transport, and limited observable peer response. No independent human semantic annotations were collected.
